\section{A fast counting algorithm via shared linear summaries}\label{sec:counting}
Sections~\ref{sec:counting}--\ref{sec:setup} prove
Theorem~\ref{thm:overview-sparse-polynomials}. We begin with the sum
from Section~\ref{subsec:preview-encoding}: for each Boolean assignment,
multiply the factors, then add the products over all assignments.
Here the factors can be any rational polynomials. We show how to
compute the sum when the replacement forms and measurements are given.
Section~\ref{sec:construction} constructs them using only the variables
read by each factor. Section~\ref{sec:setup} supplies one factor per
sparse polynomial equation.

\subsection{Polynomial factors and the value to recover}
\label{subsec:product-interface}
Let $\mathcal F$ index the factors, and let $I_f\subseteq[n]$ be the
variables read by $\Phi_f$. With rational coefficients our target is
\begin{equation}\label{eq:general-product-sum}
 Z=\sum_{x\in\{0,1\}^n}\prod_{f\in\mathcal F}\Phi_f((x_v)_{v\in I_f}).
\end{equation}
For clause satisfaction indicators, $Z$ is the number of satisfying
assignments. The proof also works with coefficients in any commutative
$\mathbb Q$-algebra $\mathcal E$: sums and products then take place in
$\mathcal E$, and interpolation uses only rational scalars. Readers may
take $\mathcal E=\mathbb Q$ throughout this section.

Different factors may use different replacements for the same variable.
Give each incidence $(f,v)$ its own indeterminate $Y_{f,v}$ and suppose
\begin{equation}\label{eq:general-product-degree}
 \deg_{\{Y_{f,v}\}}\!
    \left(\prod_{f\in\mathcal F}\Phi_f((Y_{f,v})_{v\in I_f})\right)\le d .
\end{equation}
This is the degree of the whole product computed over $\mathcal E$.
\subsection{The required rows, partitions, and measurements}
\label{subsec:layout-interface}
We use the same partitions and measurements as in Section~\ref{subsec:preview-separation},
with a group assignment for each factor in place of each clause.
The final condition below records the coefficient bound needed for
exact arithmetic. Here $\widetilde\xi$ denotes extension by zero
from a group to all $n$ coordinates.

\begin{definition}\label{def:layout}
For the factor scopes $(I_f)_{f\in\mathcal F}$, let $r\ge2$, $D\ge1$,
and $1\le R<Q$ be integers, and let $0<\tau\le1/2$.
We require the following data and properties.
\begin{enumerate}
\item For each $t\in[Q]$, a partition
$[n]=S_{t,1}\sqcup\cdots\sqcup S_{t,r}$, allowing empty parts,
with $|S_{t,g}|\le2n/r$.
\item For every factor $f\in\mathcal F$, a sequence
$\sigma_f\in[r]^Q$. The value $\sigma_f(t)$ assigns factor $f$
to a group at point $t$.
\item For every $t\in[Q]$ and $g\in[r]$, an integer matrix
$C_{t,g}\in\mathbb Z^{[c_{t,g}]\times S_{t,g}}$, where
$c_{t,g}$ is a nonnegative integer. Each row has sum of absolute
values at most $D$, and
\[
 \sum_{g=1}^r c_{t,g}\le2\tau n.
\]
For an assignment $x$, the entries of
$C_{t,g}(x|_{S_{t,g}})$ are the measurements on group $g$.
\item For every $f\in\mathcal F$ and $v\in I_f$, a row
$\vec{P}_{f,v}(T)\in\mathbb Q[T]^{1\times n}$ with
$\vec{P}_{f,v}(0)=\delta_v$ and $\deg \vec{P}_{f,v}\le R$.
For every $t\in[Q]$, $g\in[r]\setminus\{\sigma_f(t)\}$, and
$\xi\in\mathbb Q^{S_{t,g}}$, require
\[
 C_{t,g}\xi=\vec{0}\quad\Longrightarrow\quad
 \langle \vec{P}_{f,v}(t), \widetilde\xi\rangle=0,
\]
where $\widetilde\xi\in\mathbb Q^n$ is the extension of $\xi$ by zero.
\item With $d_{\mathrm{den}}=(R-1)!Q!$, every coefficient of
$d_{\mathrm{den}} \vec{P}_{f,v}$ is an integer with
$O(Q\log Q+\log D)$ bits, with an absolute implied constant.
\end{enumerate}
\end{definition}

To interpret the kernel condition, fix $t$, $f$, and $v\in I_f$,
and put $g=\sigma_f(t)$. Write $x_h=x|_{S_{t,h}}$.
If $x,x'\in\mathbb Q^n$ agree outside $S_{t,h}$ for some $h\ne g$
and satisfy $C_{t,h}x_h=C_{t,h}x'_h$, then
$x'_h-x_h\in\ker C_{t,h}$. The condition gives
\[
 \langle \vec{P}_{f,v}(t), x'\rangle=\langle \vec{P}_{f,v}(t), x\rangle.
\]
Applying this one group at a time shows that the tuple
$(\langle \vec{P}_{f,v}(t), x\rangle)_{v\in I_f}$ is determined by $x_g$ and
$(C_{t,h}x_h)_{h\ne g}$. Fixing all the measurement vectors therefore
makes the entire factor a function of $x_g$. This applies to every
polynomial factor $\Phi_f$, including factors with coefficients in
$\mathcal E$.

Each measurement of a Boolean vector lies in $[-D,D]\cap\mathbb Z$.
We bound the sum of absolute values of the matrix entries because
the construction can place the same Boolean coordinate in several
auxiliary positions. After those positions are identified, entries
of $C_{t,g}$ need not belong to $\{0,1,-1\}$. A rank bound alone
would not control the number of possible measurements: the single
row $(1,2,4,\ldots,2^{n-1})$ distinguishes all $2^n$ Boolean vectors.

\begin{proposition}[Common counting routine]\label{prop:general-count}
Let $\mathcal E$ be a commutative $\mathbb Q$-algebra. Suppose the
product of the factors in \eqref{eq:general-product-sum}, computed
over $\mathcal E$, satisfies the degree bound \eqref{eq:general-product-degree},
and that the replacement rows, partitions, and measurements satisfy
Definition~\ref{def:layout}. Put
\begin{equation}\label{eq:hasse-order}
 q=\left\lfloor\frac{dR}{Q}\right\rfloor,\qquad J=q+1.
\end{equation}
Then $Z$ can be computed by examining at most
\[
 2Qr\left(\sum_{j=0}^{\min(q,n)}2^j\binom nj\right)(2D+1)^{2\tau n}2^{2n/r}
\]
Boolean vectors. For each vector, the algorithm evaluates at most $|\mathcal F|$
factors in $\mathcal E[z]/(z^J)$. The additional operations per
vector and the final interpolation take polynomially many operations
in $\mathcal E$ and over $\mathbb Q$, in the parameters
$n,|\mathcal F|,D,Q,R,J,r$ and the total scope length.
Apart from the factor evaluators' own storage, the routine retains
only polynomially many algebra elements, rational numbers and index bits.
\end{proposition}
The proof has three parts. Lemma~\ref{lem:measurement-separation}
evaluates a sum after fixing the measurements. Lemma~\ref{lem:boolean-differences}
expresses each Taylor coefficient as restricted sums of that form.
Hermite interpolation over $\mathbb Q$ then recovers the value at zero.
Sections~\ref{subsec:measurement-factorization}--\ref{subsec:three-costs}
give these steps and their cost.

For the proof, use the supplied replacement rows, their degree
bound $R$, and the equalities
$\vec P_{f,v}(0)=\delta_v$ to form
\begin{equation}\label{eq:general-count-polynomial}
 H(T)=\sum_{x\in\{0,1\}^n}\prod_{f\in\mathcal F}
   \Phi_f\bigl((\langle\vec P_{f,v}(T),x\rangle)_{v\in I_f}\bigr)
   \in\mathcal E[T].
\end{equation}
The replacements restore the inputs at zero and multiply the degree by
at most $R$. Thus $H(0)=Z$ and $\deg_T H\le d_H:=dR$. Our application
will have $d=O(n)$, as in the bounded-occurrence example.
This uses the degree of the product over $\mathcal E$ in
\eqref{eq:general-product-degree}. Substitution is an
$\mathcal E$-algebra homomorphism: expand that product in the
independent incidence inputs, then replace each input by its
degree-at-most-$R$ form. Each surviving degree-$j$ monomial acquires
$T$-degree at most $jR$. In particular, coefficients already zero
in a truncated coefficient algebra remain zero after substitution.

\subsection{Separating the sum by fixing measurements}
\label{subsec:measurement-factorization}
Fixing the measurements makes every factor depend on only its assigned
group. The following lemma states the resulting calculation over a
commutative ring, so it applies equally to rational-valued factors and
to factors in $\mathcal E$. The restrictions $x_i=1$ will be needed
for Taylor coefficients and are local to individual groups.

\begin{lemma}\label{lem:measurement-separation}
Let $[n]=S_1\sqcup\cdots\sqcup S_r$ be a partition, allowing empty parts.
For each $g\in[r]$, let $C_g$ be an integer matrix with $c_g$ rows
and columns indexed by $S_g$. Suppose each row has sum of absolute
values at most an integer $D\ge0$, and write $c=\sum_g c_g$.
For each $g\in[r]$, let $A_g:\mathbb Q^n\to\mathcal E$, where
$\mathcal E$ is a commutative ring with identity. Suppose
\[
 A_g(x+\xi)=A_g(x)
 \qquad
 (x\in\mathbb Q^n,\ \xi\in\ker C_h,\ h\ne g),
\]
where $\xi$ is extended by zero outside $S_h$.
Then, for any $E\subseteq[n]$, the sum
\[
 Z_E=\sum_{\substack{x\in\{0,1\}^n\\x_i=1\ (i\in E)}}
           \prod_{g=1}^r A_g(x)
\]
can be evaluated by examining at most
\[
 2(2D+1)^c\sum_{g=1}^r2^{|S_g|}
\]
Boolean vectors, with at most one $A_g$ evaluation,
$\poly(n,r,c,\log(D+2))$ additional bit operations, and $O(1)$ ring
operations per vector. Apart from the evaluators' own storage and
$O(1)$ ring elements, the procedure uses $\poly(n,r,c,\log(D+2))$
bits of space.
\end{lemma}
\begin{proof}
For $y_g\in\mathbb Z^{c_g}$, define
\[
 \mathcal X_g(y_g)=\{u\in\{0,1\}^{S_g}:C_gu=y_g\},
 \qquad
 \mathcal Y=\prod_{g=1}^r\{-D,\ldots,D\}^{c_g}.
\]
Every tuple of measurements of a Boolean assignment belongs to
$\mathcal Y$, and $|\mathcal Y|=(2D+1)^c$.
We first identify the contribution of a fixed measurement tuple
$\vec y=(y_1,\ldots,y_r)$. If every $\mathcal X_g(y_g)$ is nonempty,
let $\bar x_g\in\mathcal X_g(y_g)$ be a representative for each $g$,
and put
\[
 \mathcal X_{g,E}(y_g)=
 \{u\in\mathcal X_g(y_g):u_i=1\text{ for all }i\in E\cap S_g\}.
\]
If $x_h\in\mathcal X_h(y_h)$ for every $h$, then
$x_h-\bar x_h\in\ker C_h$. Applying the invariance hypothesis
successively for $h\ne g$ gives
\[
 A_g\Bigl(\sum_{h=1}^r x_h\Bigr)
 =A_g\Bigl(x_g+\sum_{h\ne g}\bar x_h\Bigr).
\]
The assignments with measurement tuple $\vec y$ and fixed bits indexed
by $E$ form $\prod_g\mathcal X_{g,E}(y_g)$. Define
\begin{equation}\label{eq:group-sum}
 \Sigma_g(\vec y)=\sum_{x_g\in\mathcal X_{g,E}(y_g)}
       A_g\Bigl(x_g+\sum_{h\ne g}\bar x_h\Bigr).
\end{equation}
Distributivity shows that these assignments contribute
$\prod_g\Sigma_g(\vec y)$ to $Z_E$. If $\mathcal X_{g,E}(y_g)$ is empty,
then $\Sigma_g(\vec y)=0$. The representatives $\bar x_g$ themselves
need not satisfy the fixed-bit conditions indexed by $E$.
Every Boolean assignment has a unique measurement tuple, so
\[
 Z_E=\sum_{\vec y\in\mathcal Y}\prod_{g=1}^r\Sigma_g(\vec y),
\]
where the product is taken to be zero if any $\mathcal X_g(y_g)$
is empty.

The algorithm evaluates this identity one measurement tuple at a
time, starting with $Z=0_{\mathcal E}$. For each $\vec y\in\mathcal Y$,
a first scan of each $\{0,1\}^{S_g}$ finds a representative
$\bar x_g\in\mathcal X_g(y_g)$; if any such set is empty, the
algorithm skips this tuple. Otherwise it initializes
$p=1_{\mathcal E}$. To compute $\Sigma_g(\vec y)$, it initializes
$s=0_{\mathcal E}$ and scans $u\in\{0,1\}^{S_g}$ again. Whenever
\[
 C_gu=y_g,\qquad u_i=1\quad(i\in E\cap S_g),
\]
the algorithm performs the update
\[
 s\leftarrow s+A_g\Bigl(u+\sum_{h\ne g}\bar x_h\Bigr).
\]
At the end of this scan, $s=\Sigma_g(\vec y)$, and the algorithm sets
$p\leftarrow ps$. After all groups, it sets $Z\leftarrow Z+p$.
After all tuples $\vec y$ have been processed, $Z=Z_E$.

For each $\vec y$, the scan finding the representatives and the scan
computing the $\Sigma_g(\vec y)$ each examine at most
$\sum_g2^{|S_g|}$ Boolean vectors. This gives the stated bound.
Checking a vector's measurements and fixed bits uses
$\poly(n,r,c,\log(D+2))$ bit operations: each matrix row has sum of
absolute values at most $D$, so its partial sums have
$O(\log(D+2))$ bits. There is at most one factor evaluation per vector.
The retained state is $\vec y$, $(\bar x_1,\ldots,\bar x_r)$, $g$, $u$,
$s$, $p$, and $Z$. The tuple $\vec y$ contains $c$ integers in $[-D,D]$,
the representatives contain $n$ bits in total, and $s,p,Z\in\mathcal E$.
\end{proof}

To evaluate \eqref{eq:general-count-polynomial} at a fixed $t$, group
the transformed factors by $\sigma_f(t)$ and apply the lemma in
$\mathcal E$ with $E=\varnothing$. Definition~\ref{def:layout}
gives its invariance hypothesis, so the sum is $H(t)$. As in the overview, the two exponential costs
are the measurement choices and the Boolean assignments in one group.

\subsection{Taylor coefficients from fixed Boolean substitutions}
\label{subsec:fixed-boolean-substitutions}
The $Q$ values $H(1),\ldots,H(Q)$ need not determine $H$.
With $q,J$ from \eqref{eq:hasse-order}, we have $q=\lfloor d_H/Q\rfloor$.
Since $QJ>d_H$, it suffices to compute
\begin{equation}\label{eq:taylor-truncation}
 \mathcal H_t(z)=\sum_{j=0}^{J-1}([z^j]H(t+z))z^j
 \qquad(t\in[Q]).
\end{equation}
Hermite interpolation recovers $H$, and hence $Z=H(0)$, from these
$QJ$ coefficients in $\mathcal E$. Indeed, the rational Hermite
matrix for polynomials of degree less than $QJ$ is invertible over
$\mathbb Q$. Its inverse acts on elements of $\mathcal E$ by rational
linear combinations.

Fix $t$. The replacement forms at $t+z$ need not satisfy the kernel condition
at $t$. Write
\[
 \Delta\vec P_{f,v}(z)=\vec P_{f,v}(t+z)-\vec P_{f,v}(t).
\]
Every entry is divisible by $z$. For fixed $x\in\mathbb Q^n$,
introduce independent indeterminates $\vec X$ and replace each input by
\begin{equation}\label{eq:taylor-substitution}
 \langle\vec P_{f,v}(t),x\rangle
       +\langle\Delta\vec P_{f,v}(z),\vec X\rangle.
\end{equation}
Let
\begin{equation}\label{eq:auxiliary-polynomial}
 \mathcal G_x(\vec X,z)=\prod_{f\in\mathcal F}
 \Phi_f\!\left(
  (\langle\vec P_{f,v}(t),x\rangle+
   \langle\Delta\vec P_{f,v}(z),\vec X\rangle)_{v\in I_f}\right).
\end{equation}
Setting $\vec X=x$ gives
$\sum_x\mathcal G_x(x,z)=H(t+z)$.
Holding $\vec X$ fixed instead makes every correction independent
of $x$, so the kernel condition at $t$ still separates the factors.
Moreover, each occurrence of $\vec X$ carries a factor of $z$, giving
\begin{equation}\label{eq:degree-cost}
 \deg_{\vec X}([z^j]\mathcal G_x(\vec X,z))\le j.
\end{equation}
The next identity recovers the substitution $\vec X=x$ from fixed
Boolean substitutions and restrictions on the assignments being summed.

\begin{lemma}\label{lem:boolean-differences}
Let $q\ge0$ be an integer and let $\mathcal E$ be a commutative
ring with identity. For each $x\in\{0,1\}^n$, let
$p_x(X_1,\ldots,X_n)$ be a polynomial over $\mathcal E$ of degree
at most $q$. Then
\begin{align}
 \sum_{x\in\{0,1\}^n}p_x(x)
 &=\sum_{\substack{E\subseteq[n]\\|E|\le q}}
   \sum_{U\subseteq E}(-1)^{|E|-|U|}
   \sum_{\substack{x\in\{0,1\}^n\\x_i=1\ (i\in E)}}
                  p_x(\mathbf1_U).
 \label{eq:boolean-differences}
\end{align}
Here $\mathbf1_U\in\{0,1\}^n$ is the indicator vector of $U$.
\end{lemma}
\begin{proof}
For fixed $x$, we replace every positive power of $X_i$ by $X_i$.
This leaves Boolean evaluations unchanged and does not increase
the degree. By M\"obius inversion on subsets, the resulting
multilinear polynomial has the expansion
\[
 \widetilde p_x(\vec{X})=\sum_{E\subseteq[n]}c_{x,E}\prod_{i\in E}X_i,
 \qquad
 c_{x,E}=\sum_{U\subseteq E}(-1)^{|E|-|U|}p_x(\mathbf1_U).
\]
The degree bound gives $c_{x,E}=0$ whenever $|E|>q$.
The inversion is valid over the stated ring because the coefficient
of any original multilinear monomial indexed by $A$ in the
alternating sum for $c_{x,E}$ is zero unless $A\subseteq E$, and
then is
\[
 \sum_{A\subseteq U\subseteq E}(-1)^{|E|-|U|}
 =\begin{cases}1,&A=E,\\0,&A\ne E.\end{cases}
\]
Only addition and multiplication by $1$ and $-1$ are involved.
Setting $\vec{X}=x$ and summing over $x$ now gives the claimed identity:
the monomial $\prod_{i\in E}x_i$
selects exactly the assignments with $x_i=1$ for every $i\in E$.
\end{proof}

Apply Lemma~\ref{lem:boolean-differences} to
$[z^j]\mathcal G_x$ over $\mathcal E$ for each $j<J$.
All coefficients can use the same pairs $U\subseteq E$, $|E|\le q$,
so compute them together in
\[
 \mathcal R=\mathcal E[z]/(z^J).
\]
For a fixed pair $(E,U)$, let $\Phi_{f,U}(x;z)$ be factor $f$
in \eqref{eq:auxiliary-polynomial} after setting $\vec X=\mathbf1_U$.
Group these factors by their assignments:
\begin{equation}\label{eq:group-polynomial}
 A_{g,U}(x;z)=\prod_{f:\sigma_f(t)=g}\Phi_{f,U}(x;z).
\end{equation}
Their product is $\mathcal G_x(\mathbf1_U,z)$.
For $h\ne g$, changing $x$ by a vector in $\ker C_{t,h}$ changes
neither the original inputs of these factors nor their fixed
corrections. Lemma~\ref{lem:measurement-separation} therefore
computes the restricted sum with $x_i=1$ for $i\in E$ by examining
at most
\begin{equation}\label{eq:boolean-enumeration}
 2r(2D+1)^{2\tau n}2^{2n/r}
\end{equation}
Boolean vectors. Combining all the pairs gives
\begin{equation}\label{eq:full-taylor-sum}
 H(t+z)\equiv
 \sum_{\substack{E\subseteq[n]\\|E|\le q}}
 \sum_{U\subseteq E}(-1)^{|E|-|U|}
 \sum_{\substack{x\in\{0,1\}^n\\x_i=1\ (i\in E)}}
       \mathcal G_x(\mathbf1_U,z)\pmod{z^J}.
\end{equation}
The group products and sums are formed in $\mathcal R$. The result is
$\mathcal H_t(z)$; repeating at all $t$ supplies the interpolation data.
The algorithm never stores a polynomial in all $n$ indeterminates:
it sets $\vec X=\mathbf1_U$ first and stores only the $J$ coefficients in $\mathcal E$.

\subsection{The three exponential costs}
\label{subsec:three-costs}
For each $t\in[Q]$,
the number of pairs $U\subseteq E\subseteq[n]$ with $|E|\le q$ is
\begin{equation}\label{eq:pair-count}
 C_q=\sum_{s=0}^{\min(q,n)}2^s\binom ns.
\end{equation}
In particular, $C_0=1$ and $C_q=3^n$ for $q\ge n$.
For every $0<\eta\le1$,
\begin{equation}\label{eq:finite-difference-cost}
 \eta^qC_q\le\sum_{s=0}^n\binom ns(2\eta)^s=(1+2\eta)^n,
\end{equation}
because $\eta^q\le\eta^s$ for every retained index $s\le q$.

\begin{proof}[Completion of the proof of Proposition~\ref{prop:general-count}]
There are $Q$ points, $C_q$ pairs per point, and at most
\eqref{eq:boolean-enumeration} Boolean vectors per pair.
Forming each group product uses at most $|\mathcal F|$ factor evaluations.
Addition and multiplication modulo $z^J$ cost $O(J^2)$ operations
on coefficients in $\mathcal E$. The preceding Taylor identity
and interpolation prove correctness.
Process the points, pairs, measurement tuples, and local assignments
sequentially, retaining their current indices and accumulators.
Together with the $QJ$ interpolation coefficients, this gives the
stated storage bound. Bit complexity additionally depends on the
cost of evaluating the factors and on coefficient sizes; these are
verified for the sparse polynomial application in
Section~\ref{subsec:one-query-cost}.
\end{proof}

When $d=O(n)$, we have $d_H=O(nR)$ and $q=O(nR/Q)$.
Taking $\eta=R/Q$ in \eqref{eq:finite-difference-cost} gives
\[
 \log C_q\le q\log(Q/R)+n\log(1+2R/Q)
     =O\!\left(n\frac RQ\log\frac{2Q}{R}\right).
\]
Consequently, if factor evaluation and coefficient arithmetic have
polynomial bit cost, the running time is a polynomial factor times
\begin{equation}\label{eq:three-costs}
 2^{O\left(n\left[
       \frac RQ\log\frac{2Q}{R}
       +\tau\log(2D+1)+\frac1r\right]\right)}.
\end{equation}
These are exactly the three costs from the overview. The parameter
choice in Lemma~\ref{lem:incidence-parameters} makes their sum
$O((\log\log n)^2/\sqrt{\log n})$ when factor scopes and their
maximum incidence are bounded by a fixed power of $\log n$.
The routine also applies when $d$ is larger; then one keeps the
actual value $q=\lfloor dR/Q\rfloor$ in the bound of
Proposition~\ref{prop:general-count}.


